\chapter{Conjunto de instrucciones de C-Horizon}
\label{cap:instrucciones}



\section{Asignación}
Nuestro primer tipo de instrucciones son las de asignación. Como hemos visto a lo largo del documento, hay una sintaxis variada para ello dependiendo del tipo de la variable. Se muestran ejemplos de todo el documento.

\begin{algorithm}

	\Integer{} X  := \const{2};
	\comm{\\asignación a variable de tipo entero} \\
	
	\Integer{*} p  := \const{nullptr}\;
	
	p := malloc(sizeof(\Integer))
	
	p* := \const{2}
	
	\Boolean{[\const{3}]} Y \Assign 
	[\const{True}, \const{True}, \const{False}] \;
	
	\Integer{[\const{2}][\const{2}]} Z \Assign 
	[[\const{0},\const{1}],[\const{2},\const{3}]] \;
	
	\Class \texttt{estudiante} \texttt{antonio} := 
	\texttt{estudiante}(\const{1}, \const{2})\; 
	
	int(int, int) potencia(a,b) := \{ \\
	\quad \Ifff{b < \const{0}}\\
	\quad \quad \reserved{return} \const{0}\;
	\quad \} \Elsif{b == \const{0}}\\
	\quad \quad \reserved{return} \const{1}\;
	\quad \} \Elsee\\
	\quad \quad \reserved{return} potencia(a,b-\const{1})\;
	\quad \} \\
	\}\;
	
	
\end{algorithm}

Solo nos queda por comentar las funciones malloc y sizeof. La primera sirve para asignar memoria dinámica libre del tamaño específicado mientras que sizeof te dice el tamaño en palabras de 4 bytes del tipo que le das como argumento.

\section{;}
Como en C++, el ; sirve para unir dos instrucciones en una sola.

\section{While}
El siguiente tipo de instrucciones son las formadas por un while. Se usa la palabra reservada \reserved{while} y la sintaxis es la habitual,

while (expresión booleana) \{ instrucción \}
y además, obligamos a colocar \{ y \} para mostrar su ámbito. 

\section{For}
El siguiente tipo de instrucciones son las formadas por un for. Se usa la palabra reservada \reserved{for} y la sintaxis es la habitual, aunque la variable contadora se declara fuera de este.
Como antes obligamos a colocar \{ y \} para describir su ámbito. 

	
\section{If, elsif y else}
Las instrucciones if, elsif y else con la sintaxis habitual que usan palabras reservadas \reserved{if}, \reserved{elsif} y 
\reserved{else} respectivamente. Igual que en todos los casos anteriores, es obligatorio el uso de \{ \} `para declarar el ámbito de cada uno de los bloques.

\section{Switch}
Finalmente, llegamos al último tipo de instrucción, las formadas por un switch. Usamos las palabras reservadas \reserved{switch},
\reserved{case} y \reserved{default}. No es necesario usar break ya que cada switch es independiente del otro pero sí es obligatorio colocar una cláusula default al final. La instrucción switch valora una expresión booleana una sola vez y busca entre cada uno de los casos el que se corresponda con esa evaluación.

Finalmente, mostramos un ejemplo en el que se usan todas las instrucciones para ejemplificar su uso:

\begin{algorithm}

	\Integer{} X  := \const{100}\;
	\Integer{} Y  := \const{0}\;
	\Integer{} Z  := \const{1}\;
	\Integer{} T \;
	
	\reserved{int}(\reserved{int}, \reserved{int}) f(a,b) := \{ \\
	\quad \Ifff{\const{2}*a < \, b}\\
	\quad \quad \reserved{return} \const{2}*a + \const{1}\;
	\quad \} \Elsif{a > \, b}\\
	\quad \quad \reserved{return} b\;
	\quad \} \Elsee\\
	\quad \quad \reserved{return} a + \const{1}\;
	\quad \} \\
	\}\;

    \reserved{for} (Y := \const{0}; Y < \, X; Y := f(Y,X))	 \{ \\
    \quad \reserved{for} 
    (T := \const{0}; T < \, Y; T := T + \const{1})	 \{ \\
	\quad \quad \Ifff{\const{2}*Y < \, X}\\
	\quad \quad \quad Z := Z + \const{1}\;
	\quad \quad \} \\
	\quad \} \\
	\}\\
	
	
	\Integer{} Y  := X - Z\;
	
	\reserved{while} (Y < \, X)	 \{ \\
	\quad \reserved{switch}(T - \const{2}*Y)\\
	\quad \quad \reserved{case} 0 \{ \\
	\quad \quad \quad T := X\;
	\quad \quad\} \reserved{case} 1 \{ \\
	\quad \quad \quad T := 2*X\;
	\quad \quad\} \reserved{case} 2 \{ \\
	\quad \quad \quad T := \const{0}\;
	\quad \quad\} \reserved{default} \{ \\
	\quad \quad \quad T := Y\;
	\quad \quad\}\\
	\quad \} \\
	\quad Y := Y + \const{1}\;
	\}\\
\end{algorithm}

	
